% s-222-p222-m001; original PDF 222, 223
\Aphorism{s-222-p222-m001}{315}
Ist es wirklich ein Fortschritt, die Illusion zu verdünnen und einzuschränken, wenn man doch ganz ohne Illusion nicht zum Ziele kommt? Eine Frage an die realistischen Bühnendichter. Was hilft es, sich der Naturwahrheit einen Schritt mehr zu nähern als die Vorgänger, da doch die volle Naturwahrheit im Theater unmöglich ist, unmöglich als Zeit, als Raum, als Causalität — denn alle diese Grundformen des Geschehens müssen sich den Bühnenbedingungen anpassen. Das Drama fordert Umstellung der Perspectiven, Verkürzung und Verlängerung aller natürlichen Dimensionen, Anordnung des Chaotischen, Ausstreichen des Nichtssäglichen, Unterstreichen des Bedeutungsvollen — kurz eine durchweg antinaturalistische Praxis, in der sich eine naturalistische Theorie sehr als verlorener Posten ausnimmt. Es ist bereits Widernatur a priori, dass der Zuschauer etwas Rundes, Begrenztes, Stetiges zu sehen verlangt, wo das Leben nur Fragmente und Fetzen liefert. Die Wirklichkeit kennt keine geschlossenen, vor- und rückwärts beziehungsvollen Ereignisscomplexe, keine eingerahmten Bilder; sie bringt heute die „Schuld“, in zehn Jahren oder nie die „Sühne“, sie trennt Grund und Folge, Erwartung und Erfüllung um ganze Abgründe von Orts- und Zeitdifferenz, sie lässt den Blüthensamen vom Winde übers Meer tragen und in Africa statt in Europa Frucht treiben. Dieses Fehlerspiel mit unbegrenzter Dispersion, diese zerstreuende zufällige planlose Wirklichkeit nachahmen heisst für den Dramatiker auf das Drama verzichten; und \Sperr{ein wenig} nachahmen, wo völlige Nachahmung versagt bleibt, heisst mit dem Unmöglichen coquettiren. Will man das wirkliche Leben, so verbanne man von der Bühne nicht nur den Monolog, den Reim, die scandirte Rede, sondern die Rede überhaupt, jede Art Gehen und Stehen, die vorbereiteten Begegnungen, die zu Stunden zusammengedrängten Tage und Jahre; man exponire nicht, sondern überlasse dem Zuschauer, aus hingeworfenen Alltäglichkeiten klug zu werden, so wie man auch im Leben das Zusammengehörige sich erst noch zusammensuchen muss. Man lasse die Menschen auf der Bühne essen, trinken, schlafen gehen und alle vierzehn Tage ein gescheidtes Wort reden; man bringe das Durcheinander verfehlter Rendezvous, abgebrochener Pointen, gestörter Aussprachen, die ganze gähnende Breite und Nichtssäglichkeit des „wirklichen“ Lebens — man habe, wenn man natürlich sein will, den Muth zur natürlichen Stupidität!
\par

